\documentclass[11pt]{article}

\usepackage[margin=1in]{geometry}
\usepackage{amsmath, amssymb, array, booktabs, longtable}
\usepackage{natbib}
\usepackage{hyperref}

\title{Methodological Notes for Firm-Level Return Processing and CARDI Predictability Tests}
\author{CARDI Project}
\date{May 20, 2026}

\begin{document}
\maketitle

\section{Purpose}

This note documents the empirical workflow implemented in:
\[
\texttt{Code/Firm\_Test/Process\_FirmReturns}.
\]
The workflow constructs firm-level daily and monthly return datasets, estimates
firm-level pure returns, aggregates HC/MC/LC pool-level returns and volatilities,
constructs HC-versus-LC relative volatility measures, and tests whether lagged
CARDI and related indicators predict future HC relative volatility.

The code is organized around \texttt{main.R}, with module behavior controlled by
\texttt{config.R}. Helper functions are stored in:
\[
\begin{aligned}
&\texttt{functions\_load\_data.R},\quad
\texttt{functions\_returns.R},\quad
\texttt{functions\_controls.R},\\
&\texttt{functions\_pure\_return.R},\quad
\texttt{functions\_pool\_returns.R},\quad
\texttt{functions\_relative\_volatility.R},\\
&\texttt{functions\_predictability\_CARDI.R},\quad
\texttt{functions\_utils.R}.
\end{aligned}
\]

\section{Data Inputs}

The workflow combines the following data sources:

\begin{longtable}{p{0.52\textwidth}p{0.38\textwidth}}
\toprule
Input & Role\\
\midrule
\texttt{Data/Processed/Input/HighCarbonIntens} & HC stock prices and market values\\
\texttt{Data/Processed/Input/MedCarbonIntens} & MC stock prices and market values\\
\texttt{Data/Processed/Input/LowCarbonIntens} & LC stock prices and market values\\
\texttt{Data/Processed/FamaFactors/FamaFactors\_Daily.xlsx} & Daily risk-free rate\\
\texttt{Data/Processed/FamaFactors/FamaFactors\_Monthly.xlsx} & Monthly risk-free rate\\
\texttt{Data/raw/Stocklist\_industry.xlsx} & Industry classification\\
\texttt{Data/Processed/CARDI/Month\_CARDI.xlsx} & Monthly CARDI measures\\
\texttt{Data/Processed/Macro/Month\_Macro.csv} & Monthly macro controls\\
\texttt{Data/Processed/Data\_FI\_WIND.rds} & Firm financial controls\\
\texttt{Output/Carbon\_Rank.rds} & Carbon emissions ranking\\
\texttt{Output/NewIndicators/Monthly/All\_Indicators\_Monthly.csv} & Alternative indicators\\
\texttt{Data/raw/Important\_Carbon\_Events.xlsx} & Carbon-related event controls\\
\bottomrule
\end{longtable}

Each firm is assigned to a carbon group according to its source pool:
\[
\text{CarbonType}_i \in \{\mathrm{HC},\mathrm{MC},\mathrm{LC}\}.
\]

\section{Daily and Monthly Firm Returns}

\subsection{Daily Log Returns}

Let \(P_{i,d}\) denote the closing price of firm \(i\) on trading day \(d\).
The daily log return is:
\begin{equation}
R_{i,d}
=
\log P_{i,d} - \log P_{i,d-1}.
\end{equation}

Let \(r^f_d\) denote the daily risk-free rate from the daily Fama-factor file.
The daily excess return is:
\begin{equation}
XR_{i,d}
=
R_{i,d} - r^f_d.
\end{equation}

\subsection{Monthly Log Returns}

Let \(P_{i,m}^{\mathrm{end}}\) denote the price on the last trading day of month
\(m\). The monthly log return is:
\begin{equation}
R_{i,m}
=
\log P_{i,m}^{\mathrm{end}}
-
\log P_{i,m-1}^{\mathrm{end}}.
\end{equation}

Let \(r^f_m\) denote the monthly risk-free rate from the monthly Fama-factor
file. The monthly excess return is:
\begin{equation}
XR_{i,m}
=
R_{i,m} - r^f_m.
\end{equation}

Using log returns follows the standard asset-pricing convention that returns
compound additively over time.

\section{Annual Firm Controls}

\subsection{Control Variables}

The workflow constructs annual firm controls from \texttt{Data\_FI\_WIND.rds}
and related files. The core controls are:
\[
X_{i,y}
=
\left(
\log(\mathrm{Asset}_{i,y}),\;
\mathrm{ROE}_{i,y},\;
\mathrm{BM}_{i,y},\;
\mathrm{CapExRatio}_{i,y}
\right).
\]

Book-to-market is computed as:
\begin{equation}
\mathrm{BM}_{i,y}
=
\frac{\mathrm{BookValue}_{i,y}}
{\mathrm{MarketValue}_{i,y-1}^{\mathrm{last}}},
\end{equation}
where the denominator is the last observed market value from the previous year.

Capital expenditure intensity is:
\begin{equation}
\mathrm{CapExRatio}_{i,y}
=
\frac{\mathrm{CapitalExpenditure}_{i,y}}
{\mathrm{Asset}_{i,y}}.
\end{equation}

The workflow also merges lagged log carbon emissions:
\[
\log(\mathrm{Emissions}_{i,y}),
\]
but this variable is retained as an output field and is intentionally excluded
from the pure-return regression.

\subsection{Lagging Convention}

For a return observation in calendar year \(y\), the workflow uses previous-year
controls:
\begin{equation}
X^{\mathrm{lag}}_{i,t}
=
X_{i,y(t)-1}.
\end{equation}

This lagging convention reduces look-ahead bias by ensuring that controls are
measured before the return observation.

\subsection{Winsorization}

All control variables are winsorized at the 1st and 99th percentiles:
\begin{equation}
\widetilde{X}_{i,y,k}
=
\min\left\{
\max\left(X_{i,y,k}, Q_{0.01,k}\right),
Q_{0.99,k}
\right\}.
\end{equation}

Winsorization limits the influence of extreme accounting observations.

\section{Pure Return Estimation}

\subsection{Motivation}

The pure-return measure removes variation in firm excess returns associated
with observable firm characteristics and industry fixed effects. The residual is
interpreted as a return component net of these standard controls.

\subsection{Regression Specification}

For each frequency \(f \in \{\mathrm{daily}, \mathrm{monthly}\}\), the workflow
estimates a pooled OLS regression:
\begin{equation}
XR_{i,t}^{(f)}
=
\alpha
+ \beta' \widetilde{X}^{\mathrm{lag}}_{i,t}
+ \gamma_{g(i)}
+ \varepsilon_{i,t}^{(f)},
\end{equation}
where:
\[
\widetilde{X}^{\mathrm{lag}}_{i,t}
=
\left(
\mathrm{lag\_logAsset}_{i,t},
\mathrm{lag\_ROE}_{i,t},
\mathrm{lag\_BM}_{i,t},
\mathrm{lag\_CapExRatio}_{i,t}
\right),
\]
and \(\gamma_{g(i)}\) is an industry fixed effect based on the firm industry
code.

The residual is stored as:
\begin{equation}
\mathrm{PureReturn}_{i,t}^{(f)}
=
\widehat{\varepsilon}_{i,t}^{(f)}.
\end{equation}

The emission variable \(\mathrm{lag\_logEmissions}\) is not included in the
pure-return regression. This choice prevents the residualization step from
mechanically removing carbon-related variation that may be relevant for the
subsequent CARDI predictability tests.

\section{Firm-Level Monthly Volatility}

The monthly firm-level volatility measures are computed from daily observations
within each firm-month. Let \(D(i,m)\) be the set of trading days for firm \(i\)
in month \(m\). The return volatility is:
\begin{equation}
\sigma^{R}_{i,m}
=
\mathrm{sd}\left(\{R_{i,d}: d \in D(i,m)\}\right).
\end{equation}

The excess-return volatility is:
\begin{equation}
\sigma^{XR}_{i,m}
=
\mathrm{sd}\left(\{XR_{i,d}: d \in D(i,m)\}\right).
\end{equation}

The pure-return volatility is:
\begin{equation}
\sigma^{PR}_{i,m}
=
\mathrm{sd}\left(\{\mathrm{PureReturn}_{i,d}: d \in D(i,m)\}\right).
\end{equation}

In the output datasets, these variables are named:
\[
\texttt{vola\_return},\qquad
\texttt{vola\_ExReturn},\qquad
\texttt{vola\_PurReturn}.
\]

\section{Pool-Level Returns and Volatility}

\subsection{Value-Weighted Pool Returns}

For carbon group \(c \in \{\mathrm{HC},\mathrm{MC},\mathrm{LC}\}\), let
\(\mathcal{I}_{c,d}\) be the set of firms in group \(c\) with non-missing
market value on day \(d\). The value weight is:
\begin{equation}
w_{i,d}^{c}
=
\frac{\mathrm{MktCap}_{i,d}}
{\sum_{j \in \mathcal{I}_{c,d}}\mathrm{MktCap}_{j,d}}.
\end{equation}

The pool-level daily return is:
\begin{equation}
R_{c,d}
=
\sum_{i \in \mathcal{I}_{c,d}}
w_{i,d}^{c} R_{i,d}.
\end{equation}

Analogously, the workflow computes value-weighted pool excess returns and
value-weighted pool pure returns:
\begin{align}
XR_{c,d}
&=
\sum_{i \in \mathcal{I}_{c,d}} w_{i,d}^{c} XR_{i,d},\\
PR_{c,d}
&=
\sum_{i \in \mathcal{I}_{c,d}} w_{i,d}^{c} \mathrm{PureReturn}_{i,d}.
\end{align}

Value weighting follows the asset-pricing convention that larger firms receive
larger portfolio weights, as in standard portfolio construction and factor
work such as \citet{famafrench1993}.

\subsection{Pool-Level Monthly Volatility}

Pool-level monthly volatility is computed from the daily pool-level series
within each month:
\begin{align}
\sigma^R_{c,m}
&=
\mathrm{sd}\left(\{R_{c,d}: d \in m\}\right),\\
\sigma^{XR}_{c,m}
&=
\mathrm{sd}\left(\{XR_{c,d}: d \in m\}\right),\\
\sigma^{PR}_{c,m}
&=
\mathrm{sd}\left(\{PR_{c,d}: d \in m\}\right).
\end{align}

\subsection{Weighted Average Firm Volatility}

The workflow also computes a value-weighted average of firm-level monthly
volatility:
\begin{equation}
\overline{\sigma}^{R}_{c,m}
=
\sum_{i \in \mathcal{I}_{c,m}}
w_{i,m}^{c}\sigma^R_{i,m},
\end{equation}
with analogous measures for excess-return and pure-return volatility.

\section{HC-versus-LC Relative Volatility}

The main outcome variables compare HC volatility to an LC benchmark.

\subsection{Firm-Level Relative Volatility}

For an HC firm \(i\), relative volatility is:
\begin{equation}
\mathrm{Rela\_vola\_Return}_{i,m}
=
\frac{\sigma^R_{i,m}}
{\overline{\sigma}^{R}_{LC,m}}.
\end{equation}

The analogous measures are:
\begin{align}
\mathrm{Rela\_vola\_ExReturn}_{i,m}
&=
\frac{\sigma^{XR}_{i,m}}
{\overline{\sigma}^{XR}_{LC,m}},\\
\mathrm{Rela\_vola\_PurReturn}_{i,m}
&=
\frac{\sigma^{PR}_{i,m}}
{\overline{\sigma}^{PR}_{LC,m}}.
\end{align}

\subsection{Pool-Level Relative Volatility}

For the HC pool:
\begin{equation}
\mathrm{Rela\_vola\_Return}_{HC,m}
=
\frac{\overline{\sigma}^{R}_{HC,m}}
{\overline{\sigma}^{R}_{LC,m}}.
\end{equation}

The same construction is used for excess-return and pure-return volatility.

Relative volatility normalizes HC risk by the contemporaneous LC benchmark,
which helps isolate the changing risk of high-carbon firms relative to low-carbon
firms.

\section{Predictability Regressions}

\subsection{Predictor Construction}

The workflow tests whether lagged CARDI or alternative indicators predict
future HC relative volatility. For a monthly predictor \(Z_m\), the one-month
lag is:
\begin{equation}
Z_{m-1}.
\end{equation}

The main CARDI predictors include:
\[
\mathrm{CARDI\_1P\_M},\quad
\mathrm{CARDI\_5P\_M},\quad
\mathrm{CARDI\_10P\_M},
\]
and their log-difference versions:
\[
\mathrm{CARDI\_1P\_LogDiff\_M},\quad
\mathrm{CARDI\_5P\_LogDiff\_M},\quad
\mathrm{CARDI\_10P\_LogDiff\_M}.
\]

\subsection{Firm-Level Regressions}

For HC firms, the generic panel predictability regression is:
\begin{equation}
Y_{i,m}
=
\alpha_i
+ \theta Z_{m-1}
+ \delta' M_m
+ \lambda' E_m
+ \psi' C_{i,m}
+ \rho T_m
+ \eta_{y(m)}
+ u_{i,m},
\end{equation}
where:
\begin{itemize}
  \item \(Y_{i,m}\) is one of the three relative volatility outcomes;
  \item \(Z_{m-1}\) is lagged CARDI or another indicator;
  \item \(M_m\) are macro controls;
  \item \(E_m\) are event controls;
  \item \(C_{i,m}\) are firm controls;
  \item \(\alpha_i\) are firm fixed effects;
  \item \(T_m\) is a year trend;
  \item \(\eta_{y(m)}\) are year fixed effects.
\end{itemize}

The implemented firm-level specifications are:
\[
\begin{array}{ll}
\texttt{NoControls}:
& Y_{i,m} = \alpha + \theta Z_{m-1} + u_{i,m},\\[4pt]
\texttt{ControlsFirmFE}:
& Y_{i,m} = \alpha_i + \theta Z_{m-1} + \delta'M_m + \lambda'E_m + \psi'C_{i,m} + u_{i,m},\\[4pt]
\texttt{ControlsFirmFEYearTrend}:
& Y_{i,m} = \alpha_i + \theta Z_{m-1} + \delta'M_m + \lambda'E_m + \psi'C_{i,m}
  + \rho T_m + u_{i,m},\\[4pt]
\texttt{ControlsFirmFEYearFE}:
& Y_{i,m} = \alpha_i + \theta Z_{m-1} + \delta'M_m + \lambda'E_m + \psi'C_{i,m}
  + \eta_{y(m)} + u_{i,m}.
\end{array}
\]

Firm-level standard errors are clustered by firm when clustering is enabled.
This follows the finance-panel inference concerns discussed by \citet{petersen2009}.

\subsection{Pool-Level Regressions}

For HC pool-level data, the generic regression is:
\begin{equation}
Y_{m}
=
\alpha
+ \theta Z_{m-1}
+ \delta' M_m
+ \lambda' E_m
+ \rho T_m
+ \eta_{y(m)}
+ u_m.
\end{equation}

Pool-level regressions use Newey--West standard errors to account for
heteroskedasticity and serial correlation in monthly time-series residuals
\citep{neweywest1987}.

\subsection{Inference}

For time-series models, the Newey--West covariance estimator is:
\begin{equation}
\widehat{\mathrm{Var}}_{\mathrm{NW}}(\hat{\beta})
=
(X'X)^{-1}
\left(
\sum_{\ell=-L}^{L} w_\ell
\sum_{t=|\ell|+1}^{T}
x_t x_{t-|\ell|}'\hat{u}_t\hat{u}_{t-|\ell|}
\right)
(X'X)^{-1},
\end{equation}
where \(w_\ell = 1 - |\ell|/(L+1)\). The configured lag is:
\[
L = 12.
\]

For firm panels, the clustered covariance estimator allows arbitrary residual
correlation within firm:
\begin{equation}
\widehat{\mathrm{Var}}_{\mathrm{CL}}(\hat{\beta})
=
(X'X)^{-1}
\left(
\sum_{i=1}^{N} X_i'\hat{u}_i\hat{u}_i'X_i
\right)
(X'X)^{-1},
\end{equation}
with finite-sample adjustment handled by the implementation.

\section{Interpretation of the CARDI Coefficient}

The central coefficient is \(\theta\). A positive estimate implies:
\[
\theta > 0
\quad \Rightarrow \quad
\text{higher lagged CARDI predicts higher future HC relative volatility}.
\]

This is consistent with the idea that aggregate carbon-risk dynamics contain
information about future volatility pressure among high-carbon firms. The
economic motivation is related to equilibrium sustainable investing models, in
which ESG preferences affect asset demand, expected returns, realized returns
around shocks, and firms' costs of capital \citep{pastor2021}.

\section{Output Files}

\begin{longtable}{p{0.50\textwidth}p{0.40\textwidth}}
\toprule
Output & Description\\
\midrule
\texttt{Firm\_return\_daily.rds} & Firm-level daily returns, excess returns, controls, pure returns\\
\texttt{Firm\_return\_month.rds} & Firm-level monthly returns and monthly volatility\\
\texttt{HC\_daily.rds}, \texttt{MC\_daily.rds}, \texttt{LC\_daily.rds} & Pool daily returns\\
\texttt{HC\_monthly.rds}, \texttt{MC\_monthly.rds}, \texttt{LC\_monthly.rds} & Pool monthly volatility and enrichment variables\\
\texttt{HC\_Firm\_Monthly.rds} & HC firm-level relative volatility dataset\\
\texttt{HC\_Monthly.rds} & HC pool-level relative volatility dataset\\
\texttt{workflow\_validation\_summary.csv} & Observation counts and validation checks\\
\texttt{Predictability/*.csv} & Predictor-specific and model-specific regression outputs\\
\bottomrule
\end{longtable}

\section{Workflow Summary}

The complete workflow is:
\[
\begin{array}{c}
\text{Stock prices and market values}\\
\Downarrow\\
\text{Firm daily and monthly returns}\\
\Downarrow\\
\text{Lagged annual controls and industry codes}\\
\Downarrow\\
\text{Pure returns from residual regressions}\\
\Downarrow\\
\text{Firm and pool volatility measures}\\
\Downarrow\\
\text{HC-versus-LC relative volatility}\\
\Downarrow\\
\text{Lagged CARDI predictability regressions}.
\end{array}
\]

\section*{References}

\begin{thebibliography}{9}

\bibitem[Fama and French(1993)]{famafrench1993}
Fama, E. F., and French, K. R. (1993).
\newblock Common risk factors in the returns on stocks and bonds.
\newblock \emph{Journal of Financial Economics}, 33(1), 3--56.
\newblock \url{https://doi.org/10.1016/0304-405X(93)90023-5}

\bibitem[Newey and West(1987)]{neweywest1987}
Newey, W. K., and West, K. D. (1987).
\newblock A simple, positive semi-definite, heteroskedasticity and
autocorrelation consistent covariance matrix.
\newblock \emph{Econometrica}, 55(3), 703--708.
\newblock \url{https://doi.org/10.2307/1913610}

\bibitem[P\'astor, Stambaugh, and Taylor(2021)]{pastor2021}
P\'astor, L., Stambaugh, R. F., and Taylor, L. A. (2021).
\newblock Sustainable investing in equilibrium.
\newblock \emph{Journal of Financial Economics}, 142(2), 550--571.
\newblock \url{https://doi.org/10.1016/j.jfineco.2020.12.011}

\bibitem[Petersen(2009)]{petersen2009}
Petersen, M. A. (2009).
\newblock Estimating standard errors in finance panel data sets: Comparing
approaches.
\newblock \emph{The Review of Financial Studies}, 22(1), 435--480.
\newblock \url{https://doi.org/10.1093/rfs/hhn053}

\end{thebibliography}

\end{document}
